\subsection{EQUIVALENCE}

Let \(\pmb{A}\) and \(\pmb{B}\) be assemblies. The assembly

\[
(A \Longrightarrow B) \text { and } (B \Longrightarrow A)
\]

will be denoted by \(A \Longleftrightarrow B\) .

CS7. Let \(A, B, T\) be assemblies, and let \(x\) be a letter. Then the assembly \((T|x)(A \Longleftrightarrow B)\) is identical with \((T|x)A \Longleftrightarrow (T|x)B\) .

This follows immediately from CS5 (§1, no. 2) and CS6 (no. 4).

CF10. If A and B are relations in C, then \(A \Longleftrightarrow B\) is a relation in T.

This follows immediately from CF5 (§1, no. 4) and CF9 (no. 4).

If \(A \leftrightarrow B\) is a theorem in \(\mathfrak{C}\) , we shall say that \(A\) and \(B\) are equivalent in \(\mathfrak{C}\) ; if \(x\) is a letter which is not a constant of \(\mathfrak{C}\) , and if \(A\) and \(B\) are considered as relations in \(x\) , then every term in \(\mathfrak{C}\) which satisfies one also satisfies the other.

It follows from the criteria C20, C21 (no. 4) that in order to prove a theorem in \(\mathfrak{C}\) of the form \(A \Longleftrightarrow B\) , it is necessary and sufficient to be able to prove \(A \Longrightarrow B\) and \(B \Longrightarrow A\) in \(\mathfrak{C}\) . This is often done by proving \(B\) in the theory deduced from \(\mathfrak{C}\) by adjoining the axiom \(A\) , and then by proving \(A\) in the theory deduced from \(\mathfrak{C}\) by adjoining the axiom \(B\) . These remarks lead immediately to the following criteria, whose proofs we leave to the reader:

C22. Let A, B, C be relations in T. If \(A \Longleftrightarrow B\) is a theorem in T, then \(B \Longleftrightarrow A\) is a theorem in T. If \(A \Longleftrightarrow B\) and \(B \Longleftrightarrow C\) are theorems in T, then \(A \Longleftrightarrow C\) is a theorem in T.

C23. Let A, B be equivalent relations in C, and let C be a relation in C. Then the following are theorems in C:

\[
\begin{array}{c} \text {(not} A) \iff \text {(not} B); \quad (A \Rightarrow C) \iff (B \Rightarrow C); \\ (C \Rightarrow A) \iff (C \Rightarrow B); \end{array}
\]

\[
(A \text {   and   } C) \Longleftrightarrow (B \text {   and   } C); \quad (A \text {   or   } C) \Longleftrightarrow (B \text {   or   } C).
\]

C24. Let \(A, B, C\) be relations in \(\mathfrak{C}\) . Then the following are theorems in \(\mathfrak{C}\) :

(not not A) \(\Longleftrightarrow\) A; (A \(\Longrightarrow\) B) \(\Longleftrightarrow\) ((not B) \(\Longrightarrow\) (not A)); (A and A) \(\Longleftrightarrow\) A; (A and B) \(\Longleftrightarrow\) (B and A); (A and (B and C)) \(\Longleftrightarrow\) ((A and B) and C); (A or B) \(\Longleftrightarrow\) not ((not A) and (not B)); (A or A) \(\Longleftrightarrow\) A; (A or B) \(\Longleftrightarrow\) (B or A); (A or (B or C)) \(\Longleftrightarrow\) ((A or B) or C); (A and (B or C)) \(\Longleftrightarrow\) ((A and B) or (A and C)); (A or (B and C)) \(\Longleftrightarrow\) ((A or B) and (A or C)); (A and (not B)) \(\Longleftrightarrow\) not (A \(\Longrightarrow\) B); (A or B) \(\Longleftrightarrow\) ((not A) \(\Longrightarrow\) B).

C25. If A is a theorem in C and B is a relation in C, then

\[
(A \text {   and   } B) \Longleftrightarrow B
\]

is a theorem in \(\mathfrak{G}\) . If ``not \(A\)'' is a theorem in \(\mathfrak{G}\) , then \((A\) or \(B) \Longleftrightarrow B\) is a theorem in \(\mathfrak{G}\) .

¶ In principle, from now on throughout the rest of this series, criteria C1 through C25 will be used without reference.

